\begin{abstract}
Large nonlinear finite-element energy problems require reliable local
derivatives, robust nonlinear globalization, and sparse distributed linear
algebra that is matched to the mechanics of the problem. We develop a
scientific toolset for studying these requirements as coupled numerical choices
in nonlinear finite-element energy minimization and energy-based topology
optimization. The primary \jaxpetsc{} realization uses JAX for local energy and
constitutive evaluations and PETSc for sparse assembly, Newton globalization,
Krylov solvers, and preconditioners. The benchmark suite covers
$p$-Laplace, Ginzburg--Landau, finite-strain hyperelasticity, two- and
three-dimensional Mohr--Coulomb plasticity, and topology optimization.

On representative cases, we compare element automatic differentiation,
constitutive automatic differentiation, and colored sparse recovery inside PETSc
nonlinear solves. In the high-order three-dimensional plasticity
derivative-route study, constitutive AD has the lowest measured wall time among
the tested routes while matching the reported endpoint energy, external-work,
and maximum-displacement observables to displayed precision. The distributed
mechanics and topology studies then show how derivative construction interacts
with MPI assembly, nonlinear globalization, Krylov algebra, and
preconditioning. Pure JAX and FEniCS provide matched internal formulation
checks, JAX-FEM supplies a hyperelastic endpoint-state comparison, and the
Sysala-family and \v{C}erm{\'a}k--Sysala--Valdman works provide external
plasticity reference and implementation context.
Together, the evidence indicates that, for the tested problem classes, derivative
construction should be evaluated inside the sparse nonlinear solve that consumes
it, together with assembly, globalization, Krylov algebra, and preconditioning.
\end{abstract}
